\begin{pastexam}[label={px:ust-106-ee-ds-03}]{臺聯大電機類 106 學年度 第三題}
（問答）

(10\%) A bounded buffer is implemented as a circular queue. The algorithms for add and delete on the
queue are shown as follows:
\begin{minted}[fontsize=\small, style=perldoc]{c}
/* Add to a circular queue */
void add (element item)
{ if (front == (rear+1) % N)
      queueFull();
  queue[rear] = item;
  rear = (rear +1) % N;
}
/* Delete from a circular queue */
element delete()
{ element item;
  if (front == rear)
     queueEmpty();
  return queue[front];
  front = (front+1) % N;
}
\end{minted}
\begin{enumerate}
\item[(一)] (3\%) What is the capacity (maximum number of items) of the queue?
\item[(二)] (7\%) Modify both the algorithms such that the capacity of the queue is \texttt{N}.
\end{enumerate}
\end{pastexam}

\begin{pastsoln}
先說原程式兩個和題意無關的小問題，下面的改寫一併修正：
\begin{itemize}
\item \texttt{delete} 先 \texttt{return queue[front];}，下一行 \texttt{front = (front+1) \% N;}
永遠執行不到，\texttt{front} 不會前進，同一個元素會一直被取出。 應該先存進 \texttt{item}、移動
\texttt{front}，最後才 \texttt{return item;}（宣告了 \texttt{item} 卻沒用，看得出原意是這樣）。
\item \texttt{queueFull()}、\texttt{queueEmpty()} 之後程式還會往下做。 若這兩個函式只是印出錯誤
訊息，滿的時候會蓋掉資料、空的時候會讀到垃圾。 下面假設它們處理完錯誤後就 \texttt{return}。
\end{itemize}
這兩點不影響 (一) 的答案，(一) 照「\texttt{front} 有正常前進」的原意回答。

\textbf{(一)} 容量是 $N - 1$。

這裡 \texttt{front} 是前端元素的位置，\texttt{rear} 是下一個要放的位置，佇列的元素是
\texttt{queue[front]}、\texttt{queue[front+1]}…一直到 \texttt{queue[rear-1]}（索引對 $N$ 取
餘數）。 空的條件是 \texttt{front == rear}；\texttt{add} 在 \texttt{(rear+1) \% N == front} 時就
判定滿了，也就是 \texttt{rear} 再往後一格會撞到 \texttt{front} 時就不准再放。 這時佇列裡的個數是
$(\mathit{rear} - \mathit{front} + N) \bmod N = N - 1$，永遠有一格空著。

為什麼要留一格：只用 \texttt{front}、\texttt{rear} 兩個索引，固定 \texttt{front} 時 \texttt{rear}
只有 $N$ 種值，佇列的個數卻有 $0, 1, \ldots, N$ 共 $N + 1$ 種，分不出全空和全滿（都會是
\texttt{front == rear}）。 所以犧牲一格，讓「滿」變成 $N - 1$ 個（資結書第~6.2.2~節）。

例：$N = 4$，從空的開始加入 $a, b, c$ 後 $\mathit{front} = 0$、$\mathit{rear} = 3$，再加入時
$(3 + 1) \bmod 4 = 0 = \mathit{front}$，判定滿，只放得下 $3 = N - 1$ 個。

\textbf{(二)} 多用一個變數 \texttt{count} 記目前的個數（一開始 \texttt{front = rear = count = 0}）。
空是 \texttt{count == 0}，滿是 \texttt{count == N}，$N$ 格都能用。
\begin{minted}[fontsize=\small, style=perldoc]{c}
int front = 0, rear = 0, count = 0;

/* Add to a circular queue */
void add (element item)
{ if (count == N) {
      queueFull();
      return;
  }
  queue[rear] = item;
  rear = (rear + 1) % N;
  count++;
}
/* Delete from a circular queue */
element delete()
{ element item;
  if (count == 0) {
      queueEmpty();
      return ERROR_ELEMENT;   /* 自訂的錯誤值 */
  }
  item = queue[front];
  front = (front + 1) % N;
  count--;
  return item;
}
\end{minted}

\textbf{正確性}。 \texttt{count} 在每次成功的 \texttt{add} 加一、成功的 \texttt{delete} 減一，所以
一直等於佇列裡的個數；元素仍然放在 \texttt{queue[front]} 到 \texttt{queue[(front+count-1) \% N]}，
\texttt{rear} 恆等於 \texttt{(front+count) \% N}。 只在 \texttt{count == N} 時拒絕加入，所以
$N$ 格全部可用，容量是 $N$。 \texttt{front == rear} 現在可能是空也可能是滿，改由 \texttt{count}
分辨。 兩個函式仍然是 $O(1)$。

另一種做法是多記一個旗標 \texttt{tag}：\texttt{add} 後設成 $1$、\texttt{delete} 後設成 $0$；
\texttt{front == rear} 時 \texttt{tag == 1} 表示滿、\texttt{tag == 0} 表示空。 \texttt{add} 的
滿的條件改成 \texttt{front == rear \&\& tag == 1}，\texttt{delete} 的空的條件改成
\texttt{front == rear \&\& tag == 0}，容量一樣是 $N$。
\end{pastsoln}
